\providecommand{\rank}{\operatorname{r}}
\providecommand{\RR}{\mathbb R}
\providecommand{\CC}{\mathbb C}
\providecommand{\KK}{\mathbb K}

% ---------- Exercise 79 ----------
\setcounter{enumi}{78}
\item {\bfseries\boldmath
Determine a basis of each of the following vector subspaces. Do this in two ways: using the row space of a matrix and using the column space of a matrix.
}

\begin{enumerate}[label=\alph*)]

\item $V=\langle(1,2,-1),(-3,0,-3),(4,2,-2),(-1,4,-1)\rangle$ in $\RR^3$.

\begin{answer}
For the row-space method, place the generators in the rows and reduce:
the nonzero rows of an echelon form are a basis of the same row space.
For the column-space method, place the generators in the columns and
use the pivot column indices to select columns of the original matrix.

Let $v_1=(1,2,-1)$, $v_2=(-3,0,-3)$, $v_3=(4,2,-2)$,
$v_4=(-1,4,-1)$. The row matrix reduces as follows:
\[
 R=\begin{bmatrix}1&2&-1\\-3&0&-3\\4&2&-2\\-1&4&-1\end{bmatrix}
 \sim\begin{bmatrix}1&0&0\\0&1&0\\0&0&1\\0&0&0\end{bmatrix}.
\]
Thus $\bigl((1,0,0),(0,1,0),(0,0,1)\bigr)$ is a row-space basis and $V=\RR^3$.
For the column matrix $C=R^T$,
\[
 C\sim\begin{bmatrix}1&0&0&3\\0&1&0&0\\0&0&1&-1\end{bmatrix}.
\]
The pivot indices are $1,2,3$. A basis selected from the original generators
is $(v_1,v_2,v_3)$. The last column relation is $v_4=3v_1-v_3$.
\end{answer}

\item $W=\langle(2,-2,4,2),(-2,8,2,1),(-1,1,-2,-1),(0,4,8,3)\rangle$ in $\RR^4$.

\begin{answer}
Let $w_1=(2,-2,4,2)$, $w_2=(-2,8,2,1)$,
$w_3=(-1,1,-2,-1)$, $w_4=(0,4,8,3)$. For the row matrix,
\[
 R=\begin{bmatrix}2&-2&4&2\\-2&8&2&1\\-1&1&-2&-1\\0&4&8&3\end{bmatrix}
 \sim\begin{bmatrix}1&0&0&3/4\\0&1&0&1/4\\0&0&1&1/4\\0&0&0&0\end{bmatrix}.
\]
A row-space basis is
\[
 \bigl((1,0,0,3/4),(0,1,0,1/4),(0,0,1,1/4)\bigr).
\]
For $C=R^T$,
\[
 C\sim\begin{bmatrix}1&0&-1/2&0\\0&1&0&0\\0&0&0&1\\0&0&0&0\end{bmatrix}.
\]
The pivot indices are $1,2,4$, so $(w_1,w_2,w_4)$ is a second basis of $W$.
The omitted generator satisfies $w_3=-\frac12w_1$.
\end{answer}
\end{enumerate}
